\section{Introduction and statement of the result}
\label{sec:hypotheses}

The Penrose inequality relates the mass of an initial-data set to the
area needed to enclose a trapped region. Penrose's original argument
\citep{Penrose1973} connects such a bound with gravitational collapse
and the expected behavior of black holes. In the asymptotically flat
Riemannian case, Huisken and Ilmanen proved the connected-horizon
inequality by weak inverse mean curvature flow \citep{HI2001}, and
Bray proved the inequality for the total area of a possibly disconnected
horizon \citep{Bray2001}. Bray and Lee extended the Riemannian result to
dimensions below eight \citep{BrayLee2009}.

For negative cosmological constant the natural comparison metric is
Schwarzschild--anti-de Sitter. Its horizon of area $4\pi a^2$ has mass
$(a+a^3)/2$ when the cosmological constant is $-3$.
The additional cubic term is accompanied by a different asymptotic mass:
the four fluxes at spherical conformal infinity form a Lorentz covector
\citep{CH2003}. For asymptotically hyperbolic graphs,
Dahl, Gicquaud, and Sakovich established mass identities and Penrose-type
bounds \citep{DGS2013}; de Lima and Gir\~ao proved the sharp bound for
a class of balanced graphs whose horizons are star-shaped and mean
convex in the reference slice
\citep{deLimaGirao2016}. Directly passing to infinity along
inverse mean curvature flow has a distinct obstruction in this setting
\citep{Neves2010}.

Here we prove the global maximal asymptotically hyperbolic Penrose
inequality in the class specified below. Thus the corresponding maximal
formulation of the Penrose conjecture is resolved positively, with
arbitrary compact topology and without an outermostness assumption on
the original boundary. The area in the statement is the minimum
\emph{enclosing} area, not necessarily the area of that boundary.
The proof does not require a spacetime development or an additional
null-expansion condition. We make no assertion about nonmaximal
initial data or equality rigidity.

\subsection{Geometry, decay, and mass}

Let $(\Omega,g)$ be a smooth connected orientable three-manifold with
nonempty compact smooth boundary $S$. The metric space obtained by
including $S$ is complete. There is one end, with coordinates
$(r,\omega)\in(r_0,\infty)\times\Sph^2$ outside a compact region.
Write
\begin{equation}\label{eq:background}
 b=\frac{\dd r^2}{1+r^2}+r^2\sigma,
 \qquad \eta=\operatorname{arsinh}r,
 \qquad V_0=\sqrt{1+r^2},\qquad V_i=r\omega_i\quad(1\le i\le3),
\end{equation}
where $\sigma$ is the unit round metric and $\omega_i$ are the coordinate
functions of the unit sphere. The Laplacian is $\Delta=\divg\nabla$;
mean curvature is the sum of the principal curvatures.

Fix $0<\alpha<1$ and $3/2<\tau<3$. On the end, a tensor belongs to
$C^{k,\alpha}_\tau$ if its $b$-covariant derivatives through order $k$
are $O(e^{-\tau\eta})$, and its order-$k$ local H\"older seminorms
on uniform $b$-unit coordinate balls have the same weight. Equivalently,
the $C^{k,\alpha}$ norms of its components in uniform background charts,
multiplied by $e^{\tau\eta}$ at the chart centers, are bounded.
This convention is a tensor-norm convention. It does not mean
Euclidean coordinate-component decay. On compact subsets, including
the finite boundary, all fields below are smooth.

Let $K$ be a smooth symmetric covariant two-tensor on $\Omega$, smooth
up to $S$, and assume
\begin{equation}\label{eq:decay}
 g-b\in C^{2,\alpha}_\tau,
 \qquad K\in C^{1,\alpha}_\tau.
\end{equation}
For $e=g-b$, define the mass fluxes in the chosen chart by
\begin{align}\label{eq:ah-mass}
 p_\mu(g)&=\frac1{16\pi}\lim_{R\to\infty}
       \int_{\{r=R\}}\mathbb U(V_\mu,e)(\nu_b)\dd A_b,
       &&0\le\mu\le3,\\
 \mathbb U(V,e)&=V(\divg_b e-\dd\tr_b e)
              +(\tr_b e)\dd V-e(\nabla_b V,\cdot).
              \label{eq:mass-integrand}
\end{align}
Here $(\divg_b e)_i=\nabla_b^j e_{ji}$, and $\nu_b$ points toward
increasing $r$. Every operator and measure in this definition uses $b$.
We require the four limits to exist and be finite. When the covector
is future timelike, its invariant mass is
\begin{equation}\label{eq:invariant-mass}
 m_{\AH}(g)=\sqrt{p_0(g)^2-p_1(g)^2-p_2(g)^2-p_3(g)^2},
 \qquad p_0(g)>\sqrt{p_1(g)^2+p_2(g)^2+p_3(g)^2}.
\end{equation}
The spatial entries in this covector are metric mass charges; they are
not ADM momentum or angular momentum.

\subsection{Constraints, boundary, and enclosing area}

We use units $G=c=1$ and cosmological constant $\Lambda=-3$.
Maximality and the noncosmological constraint densities are
\begin{equation}\label{eq:constraints}
 \tr_gK=0,
 \qquad 16\pi\mu=R_g+6-|K|_g^2,
 \qquad 8\pi J_i=\nabla_g^jK_{ij}.
\end{equation}
Assume the dominant energy condition and weighted integrability
\begin{equation}\label{eq:dec-integrability}
 \mu\ge|J|_g,
 \qquad \int_{\mathrm{end}}V_0
       \bigl(\mu+|J|_g+|K|_g^2\bigr)\dd V_g<\infty.
\end{equation}
In particular $R_g+6\ge0$ and its $V_0$-weighted integral is finite.

On each component of $S$, choose the unit normal $\nu$ pointing into
$\Omega$, toward the designated exterior. Set
\begin{equation}\label{eq:future-trapping}
 H=\divg_S\nu,
 \qquad \tr_SK=(g^{ij}-\nu^i\nu^j)K_{ij},
 \qquad \theta_+=H+\tr_SK.
\end{equation}
We assume only $\theta_+\le0$ on every component, with this same future
orientation. For comparison with a spacetime convention, the sign is
$K(X,Y)=\overline g(\overline\nabla_Xn_{\mathrm{future}},Y)$.
There is no assumption on the inward null expansion.

\begin{definition}[Minimum enclosing area]\label{def:minimum-area}
An admissible end-side domain is a connected smooth codimension-zero
submanifold-with-boundary $D\subset\Omega$, closed as a subset of
$\Omega$, such that
\[
 \operatorname{int}D\subset\operatorname{int}\Omega
 \quad\hbox{and}\quad
 \{r>R\}\subset D\quad\hbox{for some }R.
\]
Its entire intrinsic manifold boundary $\partial D$ must be compact,
smooth, embedded, and two-sided, with any portions coinciding with $S$
included. Define
\begin{equation}\label{eq:minimum-area}
 A_{\min}(S;g)=\inf_D\Area_g(\partial D).
\end{equation}
In particular $D=\Omega$ is admissible and contributes $\partial D=S$.
An empty relative frontier is not an admissible replacement for this
intrinsic boundary.
\end{definition}

Neither the number nor the genus of the boundary components is
restricted. No trapping condition is imposed on the competitors in
Definition~\ref{def:minimum-area}.

\begin{theorem}[Global maximal AH Penrose inequality]\label{thm:main}
Under \eqref{eq:decay}--\eqref{eq:future-trapping}, the completeness,
one-end, and smoothness assumptions above, and the future-timelike
condition \eqref{eq:invariant-mass}, one has
\begin{equation}\label{eq:main-inequality}
 m_{\AH}(g)\ge\frac12\bigl(r_A+r_A^3\bigr),
 \qquad r_A=\sqrt{\frac{A_{\min}(S;g)}{4\pi}}.
\end{equation}
The coefficient is sharp: equality holds for every positive-mass
time-symmetric Schwarzschild--anti-de Sitter exterior bounded by its
minimal horizon.
\end{theorem}

The proof establishes a time-component statement for Riemannian data
which is useful independently and which does not require a causal
condition on their mass covector.

\begin{theorem}[Time-symmetric time-component inequality]\label{thm:ts}
Let $(\Omega,g)$ have the smoothness, topology, completeness, and end
structure above. Suppose $g-b\in C^{2,\alpha}_\tau$, $3/2<\tau<3$,
the four fluxes \eqref{eq:ah-mass} are finite, and
\[
 R_g\ge-6,\qquad H_g(S)\le0,
 \qquad \int_{\mathrm{end}}V_0(R_g+6)\dd V_g<\infty.
\]
In every such fixed end chart,
\begin{equation}\label{eq:ts-inequality}
 p_0(g)\ge\frac12(a+a^3),
 \qquad a=\sqrt{A_{\min}(S;g)/(4\pi)}.
\end{equation}
\end{theorem}

\subsection{Structure and new ingredients}

The first deformation solves a finite-height graph problem together
with a positive lapse and a conformal factor. It changes maximal data
into Riemannian AH data with scalar curvature at least $-6$ and a free
minimal boundary. Its mass cost tends to zero. The conformal factor
can contract area, so pointwise comparison alone is insufficient.
A volume penalty and a finite-time weak-flow argument recover the
enclosing-area comparison. A continuous choice of inner flux, based
on all possible perimeter-minimizer contact sets, removes contact with
the original weakly trapped boundary.

For Theorem~\ref{thm:ts}, a nonlinear electrostatic potential produces a
divergence-free stress. Its algebraic constraint encodes the cubic
term $a^3/2$ in the mass budget of an AF attachment. A second graph
construction on a compact interior joined to that attachment yields
an AF comparison exterior to which the ordinary Riemannian Penrose
inequality applies. We prove the coupled transmission estimates and
continuation needed for this construction. Where a trace cutoff can
spoil the scalar-curvature sign, blowup gives uniformly controlled
metric neighborhoods; additional obstacles then put these regions
strictly inside a free minimizing enclosure.

The area comparison uses weak inverse mean curvature flow only over
a fixed interval of relative areas. Its curvature estimate depends on
a scalar lower bound and fixed topology, and accommodates disconnected
leaves. It therefore has a role different from taking an asymptotic
limit of hyperbolic Hawking mass.

The numerical argument can be seen before entering these constructions.
In the positive-area time-symmetric case, fix a metric after the
preliminary approximation and write
$A_{\min}(S;g)=4\pi a^2$. Let $E_0$ denote the comparison enclosure of
the auxiliary caps alone and $E$ the enclosure that also contains the
regions where scalar curvature may be negative. Given small positive
tolerances $\eps_m,\eps_A$, the attachment and comparison constructions,
followed by area transfer, give
\[
 \begin{aligned}
 m_o&\le p_0(g)-a^3/2+\eps_m,
 &m_{\AF}(h')&\le m_o+\eps_m,\\
 \Area_{h'}(\partial E)&\ge\Area_{h'}(\partial E_0)
                    \ge(1-\eps_A)4\pi a^2.
 \end{aligned}
\]
(Propositions~\ref{prop:af-attachment}, \ref{prop:af-comparison},
and~\ref{prop:area-transfer}).
The AF Penrose inequality applies to the exterior of $E$ and yields
\[
 p_0(g)\ge \frac{a^3}{2}
       +\frac a2\sqrt{1-\eps_A}-2\eps_m.
\]
Thus the attachment supplies the cubic term and the AF inequality the
linear term. The errors tend to zero through successive finite choices;
no convergence of the comparison metrics is required.
Theorem~\ref{thm:ah-reduction} then transfers the time-symmetric bound
to maximal data. The original mass covector is put in a rest chart
before any deformation, so its time component is the invariant mass.

Sections~\ref{sec:graph-calculus} and~\ref{sec:area-transfer} develop
the shared identities and area comparison. Section~\ref{sec:maximal-reduction}
reduces Theorem~\ref{thm:main} to Theorem~\ref{thm:ts}.
Sections~\ref{sec:stress-extension}--\ref{sec:scalar-error} construct
the AF comparison for the latter. Section~\ref{sec:completion}
assembles the inequalities and checks sharpness.
